\chapter{Despliegue a producción}

La aplicación se ejecuta como un servicio Docker en el servidor personal \emph{Bigoci}. Su despliegue se
automatiza con GitHub Actions y utiliza Docker Hub como registro de imágenes,
desde el que el servidor descarga la versión que debe ejecutar.

\section{Construcción y publicación}

Cada actualización de la rama \texttt{main} activa el flujo de GitHub Actions.
Primero se construye una imagen para Linux ARM64, basada en Python 3.13, que
incluye la aplicación y sus dependencias de ejecución. La imagen se publica
en Docker Hub como \path{serhuela/fund-nav-fetcher:latest}.

Docker Hub almacena la imagen preparada para su ejecución; el despliegue
consumirá esa imagen sin volver a construirla. La etiqueta \texttt{latest} se
actualiza con cada publicación.

\section{Actualización y ejecución del servicio}

Una vez publicada la imagen, GitHub Actions copia la configuración de Docker
Compose al servidor personal \emph{Bigoci} y accede a él mediante SSH.
Descarga la imagen y recrea el contenedor con Docker Compose.

El contenedor arranca con \path{APP_ENV=production} y ejecuta la aplicación
con un usuario sin privilegios de administrador. Docker Compose configura el
reinicio automático del servicio, salvo que se haya detenido expresamente,
y limita los registros a veinte archivos de cinco megabytes cada uno.
